% !TeX program = lualatex
% Definition and sources only; this file is not an LLM solution attempt.
\documentclass[11pt,a4paper]{article}
\usepackage[margin=25mm]{geometry}
\usepackage{fontspec}
\setmainfont{lmroman10-regular.otf}[BoldFont=lmroman10-bold.otf,
 ItalicFont=lmroman10-italic.otf,BoldItalicFont=lmroman10-bolditalic.otf]
\usepackage{amsmath,amssymb,mathtools}
\usepackage{unicode-math}
\setmathfont{latinmodern-math.otf}
\AtBeginDocument{\let\setminus\smallsetminus}
\usepackage{xurl,hyperref}
\hypersetup{colorlinks=true,urlcolor=blue}
\setcounter{secnumdepth}{0}
\setlength{\parindent}{0pt}
\setlength{\parskip}{5pt}
\setlength{\emergencystretch}{3em}
\begin{document}
\section*{Boundaries of Groups and Kleinian Groups — Problem 21}\label{statement}

\subsection{Definitions and mathematical statement}
Let a group $G$ act geometrically on two proper CAT(0) spaces $X_1,X_2$. Thus both actions are proper, cocompact and isometric, and each space satisfies Euclidean upper triangle comparison. Let $Z_i=\partial_\infty X_i$ be its visual boundary, defined by bounded-distance classes of geodesic rays with the cone topology.

A compact metric space $K$ is cell-like if it has the shape of a point. Concretely, after embedding it in a Hilbert cube, it can be contracted to a point inside every neighborhood of its image. A continuous surjection is cell-like when every point preimage is cell-like.

Must there exist a compact metrizable space $Z$ and continuous surjections
\[
 Z\xrightarrow{\ p_1\ }Z_1,\qquad Z\xrightarrow{\ p_2\ }Z_2
\]
such that both $p_1$ and $p_2$ are cell-like? This is the common-domain formulation explicitly requested in the source.

Neither $Z$ nor the two maps are required to carry compatible $G$-actions. Requiring equivariance gives a stronger problem. The assertion also does not require either map to be injective: collapsing cell-like fibers is what permits different homeomorphism types. Ordinary shape equivalence is already known in this setting, but a statement of shape equivalence alone does not construct the specific common compactum and cell-like maps requested here. The same group, rather than two unrelated groups, acts on both CAT(0) models.
\subsection{Short English statement}
Do any two CAT(0) boundaries of the same group have a common compactum mapping cell-like onto both?
\subsection{Sources}
\begin{itemize}
\item[S1] ULAM.AI and the UnsolvedMath Contributors, supplied problems.json, catalogue entry, AMR Open Problem Lists. Catalogue identity and source transcription; CC BY 4.0.
\url{https://huggingface.co/datasets/ulamai/UnsolvedMath}.
\item[S2] Kapovich - Problems on Boundaries of Groups and Kleinian Groups (2005), Problem 21, PDF page 7. 
\url{https://www.math.ucdavis.edu/~kapovich/EPR/problems.pdf}.
\end{itemize}
\end{document}
